\documentclass[10pt]{article}
\usepackage[T1]{fontenc}
\usepackage[utf8]{inputenc}
\usepackage[spanish,es-nodecimaldot,es-noquoting,es-lcroman]{babel}
\usepackage{lmodern}
\usepackage{amsmath,amssymb,bm}
\usepackage[paperwidth=297mm,paperheight=500mm,margin=12mm]{geometry}
\usepackage[most]{tcolorbox}
\usepackage{tikz}
\usetikzlibrary{arrows.meta}
\usepackage{microtype}
\pagestyle{empty}
\setlength{\parindent}{0pt}
\setlength{\parskip}{2pt}
\spanishplainpercent

\newcommand{\dd}{\,\mathrm{d}}
\newcommand{\ii}{\mathrm{i}}
\newcommand{\vect}[1]{\bm{#1}}
\newcommand{\Div}{\boldsymbol{\nabla}\!\cdot}
\newcommand{\Grad}{\boldsymbol{\nabla}}
\newtcolorbox{stepbox}[2][]{%
  enhanced, sharp corners, boxrule=0.8pt, colback=blue!2,
  colframe=blue!45!black, colbacktitle=blue!9,
  fonttitle=\bfseries\large, coltitle=blue!45!black,
  left=3mm, right=3mm, top=1.5mm, bottom=1.5mm,
  before skip=0pt, after skip=0pt, #1,
  title={#2}}
\newcommand{\flow}{%
  \begin{center}\vspace{-1mm}
  \begin{tikzpicture}
    \draw[-{Stealth[length=3mm]},line width=1.1pt,blue!60!black]
      (0,0) -- (0,-3.5mm);
  \end{tikzpicture}
  \vspace{-1mm}\end{center}}

\begin{document}
\begin{center}
  {\LARGE\bfseries Del equilibrio al desplazamiento: las ecuaciones del BEM}\\[2mm]
  {\large Cilindro sólido 2D en deformación plana, respuesta armónica}\\[1mm]
  {\small Cada bloque transforma la ecuación anterior en la siguiente; las flechas indican el flujo matemático.}
\end{center}
\vspace{2mm}
\begin{tcolorbox}[enhanced,sharp corners,colback=blue!6,colframe=blue!55!black,
  boxrule=1pt,left=3mm,right=3mm,top=1mm,bottom=1mm]
\centering\large
\(\boxed{\text{ecuación de campo}}
\ \longrightarrow\ \text{integrales en la frontera}
\ \longrightarrow\ \text{sistema lineal}
\ \longrightarrow\ \text{respuesta radial}\)
\end{tcolorbox}
\vspace{2mm}

\begin{stepbox}{01 \quad Problema físico: elastodinámica en el interior \(\Omega\)}
El cilindro es infinito en \(z\); se modela su sección circular en el plano
\((x,y)\). El desplazamiento es \(\vect u=(u_x,u_y)\), con deformación plana
\((u_z=0)\). Para la convención temporal
\(u_{\text{fís}}(t)=\operatorname{Re}\{\vect u\,e^{\ii\omega t}\}\):
\[
\boxed{\Div\vect\sigma+\rho\omega^2\vect u=\vect 0\quad\text{en }\Omega},
\qquad
\boldsymbol\varepsilon(\vect u)=\tfrac12\bigl(\Grad\vect u+
(\Grad\vect u)^{\mathsf T}\bigr),
\qquad
\vect\sigma=\lambda^*\,\operatorname{tr}(\boldsymbol\varepsilon)\,\vect I+
2\mu^*\boldsymbol\varepsilon .
\]
En el borde \(\Gamma\), la normal \(\vect n\) apunta hacia afuera. La presión
uniforme apunta hacia adentro, por eso la tracción prescrita es
\(\boxed{\vect t=\vect\sigma\vect n=-p\,\vect n}\).
\end{stepbox}
\flow

\begin{stepbox}[colback=violet!2,colframe=violet!55!black,
  colbacktitle=violet!8,coltitle=violet!45!black]
{02 \quad Material complejo: amortiguamiento y velocidades de onda}
El amortiguamiento histerético se incorpora en el módulo de corte; los módulos
y las velocidades pueden ser complejos:
\[
\mu^*=\mu_0(1+2\ii\beta),\qquad
\lambda^*=\frac{2\mu^*\nu}{1-2\nu},\qquad
c_s=\sqrt{\frac{\mu^*}{\rho}},\qquad
c_p=\sqrt{\frac{\lambda^*+2\mu^*}{\rho}}.
\]
\(\rho\): densidad; \(\nu\): razón de Poisson; \(\beta\): amortiguamiento;
\(c_s,c_p\): velocidades de ondas S y P. Estos parámetros determinan los
núcleos fundamentales que aparecen más adelante.
\end{stepbox}
\flow

\begin{stepbox}{03 \quad Identidad de Green: llevar el problema a la frontera}
Al aplicar la identidad recíproca con las soluciones fundamentales, el valor
del campo en un punto fuente \(\vect\xi\) del contorno se relaciona con
desplazamientos y tracciones en todo \(\Gamma\):
\[
\boxed{c(\vect\xi)\vect u(\vect\xi)
+\operatorname{VP}\!\int_\Gamma \vect T(\vect\xi,\vect x)\vect u(\vect x)\,\dd\Gamma
=\int_\Gamma \vect U(\vect\xi,\vect x)\vect t(\vect x)\,\dd\Gamma.}
\]
\(\vect U\) es el núcleo de desplazamiento fundamental y \(\vect T\), su núcleo
de tracción. VP indica valor principal. En un borde suave \(c=\tfrac12\vect I\);
en la implementación, el término libre y la parte estática principal se
incorporan en la diagonal de \(H^s\): no se agrega otro \(\tfrac12\vect I\).
\end{stepbox}
\flow

\begin{stepbox}[colback=violet!2,colframe=violet!55!black,
  colbacktitle=violet!8,coltitle=violet!45!black]
{04 \quad Núcleos dinámicos: cómo se propaga la influencia}
Sea \(\vect r=\vect x-\vect\xi\), \(r=\|\vect r\|\),
\(\vect e=\vect r/r\), \(\gamma=c_s/c_p\),
\(z_s=\ii\omega r/c_s\), \(z_p=\ii\omega r/c_p\). Para \(K_j\), función de
Bessel modificada de segunda especie:
\[
\psi(r)=K_0(z_s)+\frac{K_1(z_s)-\gamma K_1(z_p)}{z_s},
\qquad
\chi(r)=K_2(z_s)-\gamma^2K_2(z_p),
\]
\[
\boxed{U_{ij}(\vect x,\vect\xi)=
\frac{\psi(r)\delta_{ij}-\chi(r)e_i e_j}{2\pi\mu^*}.}
\]
El núcleo de tracción se obtiene aplicando la ley constitutiva al campo
fundamental y proyectando el esfuerzo sobre la normal exterior en
\(\vect x\):
\[
T_{ij}(\vect\xi,\vect x)=
\bigl[\boldsymbol\sigma_{\vect x}(\vect U_{i\cdot})\bigr]_{jk}n_k(\vect x).
\]
Aquí \(i\) identifica la fuerza auxiliar y \(j\), la componente de tracción:
se aplica el operador a la fila \(U_{i\cdot}\). En la identidad de frontera,
\(U\) pondera las tracciones y \(T\), los desplazamientos.
\end{stepbox}
\flow

\begin{stepbox}{05 \quad Discretizar \(\Gamma\): elementos cuadráticos e integración}
Solo se divide el contorno; no se crea una malla del interior. Con \(n_e\)
elementos hay \(N=2n_e\) nodos globales y dos componentes \((u_x,u_y)\) por nodo.
En cada elemento, \(\eta\in[-1,1]\) y los nodos locales están en
\((-1,0,1)\):
\[
N_1=\tfrac12\eta(\eta-1),\qquad N_2=1-\eta^2,\qquad
N_3=\tfrac12\eta(\eta+1),\qquad
\vect x(\eta)=\sum_{a=1}^3N_a(\eta)\vect x_a.
\]
La geometría curva aporta el jacobiano y la normal
\[
J=\left\|\frac{\dd\vect x}{\dd\eta}\right\|,\qquad
\vect n=\frac{(y',-x')}{J},\qquad
\dd\Gamma=J\,\dd\eta,\qquad
\int_{\Gamma_e}f\,\dd\Gamma\ \approx\ \sum_q f(\vect x_q)J_qw_q.
\]
Para una función de forma, el peso integrado es
\(W_{qa}=N_a(\eta_q)J_qw_q\). Cuando el punto fuente cae en el elemento, la
cuadratura se adapta alrededor de esa singularidad.
\end{stepbox}
\flow

\begin{stepbox}[colback=teal!2,colframe=teal!55!black,
  colbacktitle=teal!8,coltitle=teal!45!black]
{06 \quad Ensamblar \(H\) y \(G\): convertir integrales en bloques \(2\times2\)}
Al evaluar la identidad en cada nodo de colocación \(i\), se interpola el campo
con las funciones \(N_g\) de los nodos globales. La formulación discreta es
\[
H_{ig}=C_{ig}+\operatorname{VP}\!\int_\Gamma N_g(\vect x)\,T\,\dd\Gamma,
\qquad
G_{ig}=\int_\Gamma N_g(\vect x)\,U\,\dd\Gamma,\qquad
C_{ig}=c(\vect\xi_i)\delta_{ig}.
\]
El código separa el operador estático singular de la corrección dinámica:
\[
H_{ig}=H^s_{ig}+\int_\Gamma N_g(T-T^s)\,\dd\Gamma,\qquad
G_{ig}=\int_\Gamma N_gU\,\dd\Gamma,\qquad
H^s_{ii}=-\sum_{g\ne i}H^s_{ig}.
\]
La última igualdad hace que cada fila de bloques de \(H^s\) sume cero
(traslación rígida estática); la diagonal no se omite. Al final,
\[
\boxed{H(\omega)\vect u=G(\omega)\vect t},
\qquad H,G\in\mathbb C^{\,2N\times2N},\quad
\vect u,\vect t\in\mathbb C^{\,2N}.
\]
\end{stepbox}
\flow

\begin{stepbox}[colback=orange!3,colframe=orange!65!black,
  colbacktitle=orange!10,coltitle=orange!55!black]
{07 \quad Aplicar la carga, resolver y extraer la respuesta buscada}
La presión da el vector conocido \(\vect t=-p\vect n\); se resuelve el sistema
lineal para las dos componentes del desplazamiento complejo en cada nodo:
\[
\boxed{\vect u=\operatorname{solve}\!\left(H,G\vect t\right).}
\]
Para el círculo, en el nodo \(i\), \(\vect n_i=\vect x_i/R\). La salida radial es
\[
\boxed{u_{r,i}=\vect u_i\cdot\vect n_i
=u_{x,i}n_{x,i}+u_{y,i}n_{y,i}.}
\]
\(u_{r,i}\) es complejo: \(|u_{r,i}|\) da la amplitud local y
\(\arg(u_{r,i})\) su fase con la convención \(e^{\ii\omega t}\). La clase devuelve
un valor radial por nodo; no devuelve solo un promedio.
\end{stepbox}

\vfill
\begin{tcolorbox}[enhanced,sharp corners,colback=black!4,colframe=black!35,
  boxrule=0.6pt,left=3mm,right=3mm,top=1mm,bottom=1mm]
\centering\small
\textbf{Lectura global:}\quad
\(\Div\sigma+\rho\omega^2u=0\)
\(\xrightarrow{\text{Green}}\)
\(cu+\int_\Gamma Tu=\int_\Gamma Ut\)
\(\xrightarrow{\text{elementos}}\)
\(H(\omega)u=G(\omega)t\)
\(\xrightarrow{t=-pn}\)
\(u_r=u\cdot n\).
\end{tcolorbox}
\begin{center}\scriptsize
Esquema matemático del modelo de \texttt{BEMCilindro}; fuente de ecuaciones:
\texttt{estudio\_bignardi/bem\_cilindro.py}.
\end{center}
\end{document}
